% 有向环与有向无环图：长度为2的有向环；三角形的循环定向；三角形的无环定向1->2,1->3,2->3。
\begin{tikzpicture}[panelmark/.style={font=\figurefont\fontsize{8}{10}\selectfont,text=FigureInk,anchor=north,inner sep=0pt,outer sep=0pt},
>=Stealth,
  v/.style={circle,draw=FigureStroke,fill=FigureBlueFill,line width=.8pt,minimum size=6.5mm,inner sep=0pt},
  a/.style={draw=FigureStroke,line width=.8pt,->,shorten >=1pt},
  ar/.style={draw=FigureRed,line width=1.2pt,->,shorten >=1pt}]
\begin{scope}[local bounding box=panelrow]
\begin{scope}[local bounding box=panela]
  \node[v] (u) at (0,0) {$i$}; \node[v] (w) at (1.5,0) {$j$};
  \draw[ar] (u) to[bend left=30] (w); \draw[ar] (w) to[bend left=30] (u);

\end{scope}
\begin{scope}[xshift=3.6cm,local bounding box=panelb]
  \node[v] (1) at (0,.75) {$1$}; \node[v] (2) at (-.85,-.6) {$2$}; \node[v] (3) at (.85,-.6) {$3$};
  \draw[ar] (1)--(2); \draw[ar] (2)--(3); \draw[ar] (3)--(1);

\end{scope}
\begin{scope}[xshift=6.9cm,local bounding box=panelc]
  \node[v] (1) at (0,.75) {$1$}; \node[v] (2) at (-.85,-.6) {$2$}; \node[v] (3) at (.85,-.6) {$3$};
  \draw[a] (1)--(2); \draw[a] (2)--(3); \draw[a] (1)--(3);

\end{scope}
\end{scope}
% Subfigure marks: 8 pt, centred below each panel; common top 3 mm below this row.
\coordinate (panelrowbase) at ($(panelrow.south)+(0,-3mm)$);
\node[panelmark] at (panela.center |- panelrowbase) {(a)};
\node[panelmark] at (panelb.center |- panelrowbase) {(b)};
\node[panelmark] at (panelc.center |- panelrowbase) {(c)};
\end{tikzpicture}
